\documentclass[a4paper,11pt]{article}
\usepackage[T2A]{fontenc}
\usepackage[utf8]{inputenc}
\usepackage[russian]{babel}
\usepackage{amsmath,amssymb,mathtools}
\usepackage{helvet}
\renewcommand{\familydefault}{\sfdefault}
\usepackage{icomma}
\usepackage{geometry}
\geometry{left=24mm,right=23mm,top=23mm,bottom=25mm}
\usepackage{microtype}
\usepackage{xcolor}
\usepackage{tikz}
\usepackage{fancyhdr}
\usepackage{hyperref}
\usepackage{setspace}
\setstretch{1.16}
\definecolor{accent}{HTML}{008080}
\definecolor{darkaccent}{HTML}{004D4D}
\definecolor{ink}{HTML}{303638}
\definecolor{muted}{HTML}{5D6668}
\hypersetup{colorlinks=true,linkcolor=darkaccent,urlcolor=darkaccent,
  pdfauthor={Лёвин Артём Александрович},
  pdftitle={Проверка домашней работы — Ярослав Гаврилов — 23.11.2026}}
\setlength{\parindent}{0pt}
\setlength{\parskip}{7pt}
\setlength{\headheight}{15pt}
\pagestyle{fancy}
\fancyhf{}
\fancyhead[L]{\small\color{muted}Задачи на движение}
\fancyhead[R]{\small\color{muted}Ярослав Гаврилов}
\fancyfoot[C]{\small\color{muted}\thepage}
\renewcommand{\headrulewidth}{0.3pt}
\renewcommand{\headrule}{\hbox to\headwidth{\color{accent}\leaders\hrule height \headrulewidth\hfill}}
\newcommand{\sectlabel}[1]{\par\medskip\textbf{\color{darkaccent}#1}\par\vspace{-2pt}}
\begin{document}
\thispagestyle{empty}
\begin{center}
\vspace*{28mm}
{\large\color{darkaccent}МАТЕМАТИКА}\par
\vspace{24mm}
{\LARGE\color{ink}Проверка домашней работы}\par
\vspace{12mm}
{\Large\color{accent}Задачи на движение по реке}\par
\vspace{11mm}
{\normalsize\color{muted}Подробный разбор задания, требующего завершения}\par
\vfill
{\large Ярослав Гаврилов}\par
\vspace{10mm}
{\large 23 ноября 2026 г.}\par
\vfill
{\small Лёвин Артём Александрович}\par
{\small эксклюзивно для Ярослава Гаврилова}\par
\vspace{15mm}
\begin{tikzpicture}
\draw[accent,line width=0.9pt] (-4.8,0) .. controls (-3.5,0.65) and (-1.7,-0.65) .. (0,0)
.. controls (1.6,0.65) and (3.4,-0.65) .. (4.8,0);
\end{tikzpicture}
\end{center}
\clearpage
\section*{\color{darkaccent}Задание № 9. Движение по течению и против течения}
\vspace{-3mm}
\textcolor{accent}{\rule{\linewidth}{0.45pt}}

\sectlabel{Условие}
Судно проходит $36$ км против течения на $3$ часа дольше, чем такое же расстояние по течению. Скорость течения равна $3$ км/ч. Найдите собственную скорость судна.

\sectlabel{Что требуется исправить}
На предоставленных фотографиях решение девятого задания не представлено. Чтобы решить такую задачу, сначала надо выразить две скорости судна относительно берега, затем --- два времени и составить уравнение по разнице этих времён.

\sectlabel{Ключевая идея}
Собственная скорость судна --- это его скорость в стоячей воде. По течению течение \emph{помогает} судну, а против течения --- \emph{мешает}:
\[
\begin{aligned}
\text{скорость по течению}&=v+3,\\
\text{скорость против течения}&=v-3.
\end{aligned}
\]
Здесь $v$ --- собственная скорость судна в км/ч. Поскольку судно движется против течения, $v>3$.

\sectlabel{Подробное решение}
\textbf{Шаг 1.} Используем формулу времени движения:
\[
 t=\frac{s}{v}.
\]
Оба расстояния равны $36$ км, поэтому
\[
 t_{\text{против}}=\frac{36}{v-3},\qquad
 t_{\text{по}}=\frac{36}{v+3}.
\]

\textbf{Шаг 2.} Против течения судно идёт на $3$ часа дольше. Значит, из большего времени вычитаем меньшее:
\[
\frac{36}{v-3}-\frac{36}{v+3}=3.
\]

\textbf{Шаг 3.} Приводим левую часть к общему знаменателю:
\[
\begin{aligned}
\frac{36(v+3)-36(v-3)}{(v-3)(v+3)}&=3,\\[2pt]
\frac{36\bigl(v+3-v+3\bigr)}{v^2-9}&=3,\\[2pt]
\frac{216}{v^2-9}&=3.
\end{aligned}
\]

\textbf{Шаг 4.} Учитывая $v>3$, умножаем обе части на ненулевой знаменатель и находим скорость:
\[
\begin{aligned}
216&=3(v^2-9),\\
72&=v^2-9,\\
v^2&=81,\\
v&=9.
\end{aligned}
\]
Отрицательное значение $-9$ не подходит: скорость не может быть отрицательной, и требуется $v>3$.

\sectlabel{Проверка результата}
При собственной скорости $9$ км/ч скорость против течения равна $9-3=6$ км/ч, а по течению --- $9+3=12$ км/ч. Тогда времена движения:
\[
\frac{36}{6}=6\text{ ч},\qquad \frac{36}{12}=3\text{ ч}.
\]
Разница составляет ровно $6-3=3$ часа, как и указано в условии.

\sectlabel{Ответ}
\[\boxed{9\text{ км/ч}}\]

\sectlabel{Задача на закрепление}
Катер проходит $48$ км против течения на $2$ часа дольше, чем такое же расстояние по течению. Скорость течения равна $2$ км/ч. Найдите собственную скорость катера.

\vfill
{\footnotesize\color{muted}Лёвин Артём Александрович эксклюзивно для Ярослава Гаврилова}
\end{document}
